\documentclass[a4paper]{article}

\usepackage{verbatim}
\usepackage{graphics}
\usepackage{geometry}
\usepackage[T1]{fontenc}
%\usepackage[utf8]{inpuntenc}
\usepackage{ragged2e}
\usepackage{parskip}
\setlength{\parskip}{1em}
\usepackage{fancyhdr}
\usepackage{lastpage}
\usepackage{multicol}

\usepackage[most]{tcolorbox}
\tcbuselibrary{listingsutf8, minted}
\usepackage{minted}
\usepackage{xcolor}

% Define a caixa estilizada para código
\newtcblisting{meuCodigo}{
    listing engine=minted,
    colback=white,
    colframe=black!75!white,
    boxrule=0.5pt,
    arc=4pt,
    outer arc=4pt,
    top=6pt, bottom=6pt, left=8pt, right=8pt,
    minted options={
        %linenos=true,
        breaklines=true,
        tabsize=4,
        style=monokai,
        fontsize=\small,
    }
}

\usepackage[brazil]{babel}
\usepackage{hyperref}
\usepackage{fontspec}

\usepackage{amsfonts, amsmath, amssymb}
\usepackage{mathtools}
\usepackage{physics}
%\usepackage{siunitx}

\title{SQL - \textit{\textbf{\href{https://www.udemy.com/course/curso-completo-sql-para-analise-de-dados/}{Ûdemy}}}}

\author{Hygor Melo de Sena}

\date{07 de maio de 2025}

\setmainfont{Spectral}

\sloppy

\renewcommand{\headrulewidth}{0.5pt}
\renewcommand{\footrulewidth}{0.5pt}

\setlength{\columnseprule}{1pt}
\setlength{\columnsep}{13em}

\pagestyle{fancy}

\begin{document}
    \maketitle
        
        \newpage
        
    \tableofcontents
        
        \newpage
    
    \section{Módulo 3:}
    
    
        %\begin{minted}[fontsize=\small, breaklines]{sql}
           
        %\end{minted}
    
    
        \subsection{Comando de Seleção}
            
            Em \textit{DQL - ( Data Query Languagem )}, que é um subconjunto de \textit{SQL}, e o foco do estudo, o comando \textbf{select} se faz essencial
            
            \bigskip
            
            \textbf{select}
            
                \begin{itemize}
                    \item Único comando obrigatório;
                    \item \textbf{select} = selecionar;
                    \item Não obrigatório consultar um banco de dados;
                    \item Resultado da consulta é outra tabela; \footnote{Keep in Mind - É importante se atentar a esse detalhe.}
                    \item Você escolhe o que deseja;
                \end{itemize}
            
            \bigskip
            
            \noindent Abaixo teremos dois exemplos de como podemos usar o \textit{SQL}. Observe que estamos utilizando o \textit{BigQuery} e que criamos uma \textit{dataset} chamado \textit{ecommerce}, com o seguintes bancos de dados: \textit{categories, customers, items, orders, products.}
            
            \bigskip
            
            Exemplo 1.1 \(\lrcorner\)
                
                \begin{minted}[fontsize=\small, breaklines]{sql}
                
                # SELECT            
                SELECT
                4 as caluna_1,
                10 as coluna_2,
                "Meu nome é Caio" as nome,
                "Meu nome é Caio" as nome_2,
                5*3 as multiplicacao,
                20/4 as divisao;
                \end{minted}
        
            \bigskip
            
            Exemplo 1.2 \(\lrcorner\)
            
                \begin{minted}[fontsize=\small, breaklines]{sql}
                # SELECT                
                SELECT
                    name
                FROM `cursos-sql.ecommerce.categories`;
                
                \end{minted}
            
            No exemplo 1.1, temos elementos numéricos, strings e operações matemáticas sendo associadas a uma dada coluna de nome a escolha do programador. Perceba que um mesmo elemento pode ser associado a colunas distintas no banco de dados como vemos com o \texttt{"Meu nome é Caio" as nome} e \texttt{"Meu nome é Caio" as nome\_2}. Note também que para finalizarmos os blocos de comandos utilizamos o \textit{;} como podemos ver em \texttt{20/4 as divisao}. Além disso, saiba que o \textit{SQL} é \textit{case sensitive}, ou seja, em algumas partes do programa faz diferença em \textbf{CAIXA ALTA} ou \textbf{minusculas}; ou mesmo mesclar ambas.
            
            No exemplo 1.2, temos um chamado de um banco de dados. Observemos agora o comando \textit{FROM}. No comando \textit{FROM} tem o seguinte: \texttt{FROM `cursos-sql.ecommerce.categories`;}. Primeiramente temos os sinais de crase \texttt{`\(\ldots\)`} \hspace{0.1em}, utilizados para uma melhor escrita e relacionar diferentes \textit{datasets}. Segundo teremos o nome do projeto (onde esta todos os nossos \textit{datasets}) seguido do \textit{dataset} e do referido banco de dados do qual desejamos obter informações.
            
        \subsection{Comando Limit e Distinct}
        
            O comando \textit{LIMIT} é usado para definir a quantidade de linha que se deseja coletar de uma coluna. Veja no exemplo abaixo:
            
            \bigskip
            
            Exemplo 1.3 \(\lrcorner\)
            
            \begin{minted}{sql}
                # SELECT
                SELECT
                    * 
                FROM `cursos-sql.ecommerce.categories`
                LIMIT 5;
            \end{minted}
            
            \bigskip
            
            No exemplo 1.3 temos dois comando novos, o asterisco \textit{'*'} e o \textit{LIMIT}. O asterisco significa selecionar todas as colunas de uma determinada tabela dentro de um \textit{dataset} e o \textit{LIMIT} limita a quantidade de linha exibidas da tabela selecionada.   
            
            O comando \textit{DISCTINCT} é utilizado para mostrar quantos elementos distinto eu tenho dentro de um coluna, ou seja, será ignorado os valores repetidos. Além disse, o \textit{DISTINCT} pode ser usado em conjunto com funções de agregação. Veja abaixo:
            
            \bigskip
            
            Exemplo 1.4 \(\lrcorner\)
            
            \begin{minted}{sql}
                SELECT DISTINCT
                    category_id
                FROM ecommer.products;
                
                SELECT DISTINCT
                    first_name
                FROM ecommer.products;
                
                SELECT
                    COUNT(DISTINCT first_name) as nomes_distintos
                FROM ecommer.products;
            \end{minted}
            
            \bigskip
            
            A função de agregação \textit{COUNT} contará quantos elementos distintos estão presentes dentro da coluna \textit{first\_name} e exibirá o resultado dentro da coluna apelido nomes\_distintos. 
            
        \subsection{Operadores}
        
        Operadores são utilizados para vários fins e propósitos, vejamos alguns deles:
        
        \bigskip
        
        \begin{multicols}{2}
            Aritméticos\(\lrcorner\)
            
            \begin{itemize}
                \item \texttt{+} (soma)
                \item \texttt{-} (subtração)
                \item \texttt{*} (multiplicação)
                \item \texttt{/} (divisão)
            \end{itemize}
            
            \medskip
            
            Comparação\(\lrcorner\)
            
            \begin{itemize}
                \item \texttt{=} (igual)
                \item \texttt{<> ou !=} (diferente)
                \item \texttt{>=} (maior ou igual)
                \item \texttt{<=} (menor ou igual)
                \item \texttt{>} (maior)
                \item \texttt{<} (menor)
                \item \texttt{between} (entre)
                \item \texttt{like} (como)
                \item \texttt{in} (está contido)
            \end{itemize}
            
            \medskip
            
            IS (é no sentido de existir)\(\lrcorner\)
            
            \begin{itemize}
                \item \texttt{is null} (é nulo)
                \item \texttt{is true} (é verdadeiro)
                \item \texttt{is false} (é falso)
            \end{itemize}
            
            \medskip
            
            Lógicos\(\lrcorner\)
            
            \begin{itemize}
                \item \texttt{NOT} (não)
                \item \texttt{AND} (e)
                \item \texttt{OR} (ou)
            \end{itemize}
        \end{multicols} 
        
        \subsubsection{Aplicação dos Operadores}
        
        Abaixo temos a aplicação de operadores aritméticos:
        
        \bigskip
        
        \begin{minted}{sql}
            # aritiméticos
            SELECT
                7 + 7 as soma,
                12 - 4 as subtracao,
                73 * 4 as multiplicacao,
                1 / 3 as divisao;
        \end{minted}
        
        \bigskip
        
        \pagebreak[1-4]
        
        Aplicação dos operadores de comparação:
        
        \bigskip
        
        \begin{minted}{sql}
            # comparação
            SELECT
                5 = 5,
                1 > 10,
                100 < 3,
                60 <= 60,
                17 != 17;
                
            # between
            SELECT
                10 between 8 and 15,
                -13 between 0 and 100,
                3 between 3 and 5,
                5 between 3 and 5;
                
                
            SELECT
                *
            FROM ecommerce.products
            WHERE price between 30 and 80;
            
            # like
            SELECT
                *
            FROM ecommerce.custumers
            WHERE first_name like 'J%' and '%r';    
             
            SELECT
                *
            FROM ecommerce.custumers
            WHERE first_name like 'J%r';
            
            #in
            
            SELECT
                *
            FROM ecommerce.products
            WHERE price in (10, 11, 15);
            
            SELECT
                *
            FROM ecommerce.products
            WHERE state in ("Acre", "Ceará");             
        \end{minted}
               
        \bigskip
           
\end{document}